\documentclass[final,5p,times,authoryear]{elsarticle}

\usepackage{amsmath,amssymb,amsfonts}
\usepackage{amsthm}
\usepackage{graphicx}
\usepackage{subfig}
\usepackage{booktabs}
\usepackage{algorithm}
\usepackage{algpseudocode}
\usepackage{textcomp}
\usepackage[hidelinks]{hyperref}
\usepackage{microtype}          % better justification; removes most overfulls
\setlength{\emergencystretch}{2em} % last-resort stretch for math-heavy lines

\usepackage{placeins}           % \FloatBarrier: keep every float before the Conclusion

%% Float-placement tuning: allow the wide figure* floats to share a page
%% top with text (up to three per page) so they are placed as soon as
%% possible instead of drifting past the Conclusion.
\setcounter{topnumber}{4}
\setcounter{dbltopnumber}{3}
\setcounter{totalnumber}{8}
\renewcommand{\topfraction}{0.92}
\renewcommand{\dbltopfraction}{0.92}
\renewcommand{\textfraction}{0.07}
\renewcommand{\floatpagefraction}{0.75}
\renewcommand{\dblfloatpagefraction}{0.75}

\theoremstyle{plain}
\newtheorem{theorem}{Theorem}
\newtheorem{lemma}{Lemma}
\newtheorem{proposition}{Proposition}
\theoremstyle{definition}
\newtheorem{assumption}{Assumption}
\theoremstyle{remark}
\newtheorem{remark}{Remark}

\journal{Automatica}

%% Suppress the "Preprint submitted to Automatica <date>" first-page
%% footer (keep only a centred page number).
\makeatletter
\def\ps@pprintTitle{%
  \let\@oddhead\@empty
  \let\@evenhead\@empty
  \def\@oddfoot{\centerline{\thepage}}%
  \let\@evenfoot\@oddfoot}
\makeatother

\begin{document}

\begin{frontmatter}

\title{Model-Free Low-and-High Gain Q-Learning for Discrete-Time
Linear Systems with Actuator Saturation}

\author[ttu]{Md~Nur-A-Adam~Dony}
\author[ttu]{Syed~Ali~Asad~Rizvi\corref{cor}}
\ead{srizvi@tntech.edu}
\author[uva]{Zongli~Lin}
\cortext[cor]{Corresponding author.}
\address[ttu]{Department of Electrical and Computer Engineering,
Tennessee Technological University, Cookeville, TN 38505, USA}
\address[uva]{Charles L.~Brown Department of Electrical and Computer
Engineering, University of Virginia, Charlottesville, VA 22903, USA}

\begin{abstract}
A reinforcement-learning route to low-and-high gain control is
developed for discrete-time linear systems whose actuators saturate,
covering state as well as output feedback and requiring no knowledge
of the plant matrices.  Q-learning with a parameterized Riccati cost
supplies the low-gain policy, and the central observation is that the
converged Q-kernel already encodes---behind one Schur complement---the
associated Lyapunov matrix, from which the largest control magnitude
attainable anywhere on the Lyapunov level set is available in closed
form.  A gain boost certified jointly for the actuator limits and for
strict Lyapunov decrease, together with the exact stability margin of
that boost, is then read off from a single scalar of the kernel; no
identification step, further learning, or numerical search is
involved.  The output-feedback extension reconstructs
the state from a sliding window of past inputs and outputs, and the
entire state-feedback pipeline transports under one
information-preserving substitution: the Lyapunov-matrix inverse is
replaced by the Moore--Penrose pseudoinverse of its rank-deficient
extended counterpart.  Numerical experiments on a fourth-order
benchmark validate both formulations and reveal a structural trade-off
between gain-boost margin and pole-movement effectiveness, controlled
by the rank of the cost weighting.
\end{abstract}

\begin{keyword}
Reinforcement learning \sep Q-learning \sep actuator saturation \sep
low-and-high gain feedback \sep model-free control \sep output feedback
\end{keyword}

\end{frontmatter}

\section{Introduction}\label{sec:intro}
%% ======================================================================

Physical actuators deliver only bounded control effort, making input
saturation one of the most pervasive nonlinearities in control
practice.  For linear plants, global stabilizability by bounded inputs
is characterized by asymptotic null-controllability
(ANCBC)~\citep{SussmannSontagYang1994}, with the discrete-time theory
in~\citep{YangSontagSussmann1997}.  Semi-global designs built on
low-gain feedback~\citep{LinSaberi1993,Lin1998} keep the actuator
inside its limits by letting the gain shrink toward zero, at the price
of sluggish transients; low-and-high gain
compensation~\citep{LinSaberi1995,HuLin2001} recovers speed by
amplifying the low-gain law, and in discrete time the amplification is
necessarily
finite~\citep{LinSaberiStoorvogelMantri2000,WangGripSaberiJohansen2012}.
The discrete-time constructions available to
date~\citep{LinSaberiStoorvogelMantri2000} presuppose exact knowledge
of the plant matrices.

Data-driven synthesis removes this model dependence.  Within
reinforcement learning~\citep{SuttonBarto2018,Bertsekas2019},
Q-learning~\citep{WatkinsDayan1992} is attractive because the optimal
policy is obtained directly from trajectories, with no identification
stage~\citep{LewisVrabieVamvoudakis2012,JiangJiang2012,Kleinman1968}.
For saturated discrete-time plants, however, the model-free schemes
of~\citep{RizviLin2018CDC,RizviLin2019IJRNC,RizviLin2023book} verify
the actuator constraint only \emph{pointwise}, at states that are
already contracting; the resulting certificate is loose, most of the
available control authority is left idle, and no low-and-high gain
construction without a model has been reported
(Table~\ref{tab:related_work}).



Both restrictions are removed here, for state and for output
feedback.  The point of departure is that a Schur complement of the
converged Q-kernel reproduces the Lyapunov matrix of the low-gain
design; from it, the largest control demanded anywhere on the
invariant level set---and with it a boost bound certifying saturation
avoidance jointly with strict Lyapunov decrease---comes out in closed
form from one scalar of the kernel, without a model, further data, or
search.  In the model-free setting, the resulting stability limit
plays the role of the discrete-time constraint
of~\citep{WangSaberiStoorvogelSannuti2012}.  On the output-feedback
side we adopt the parameterization of \citet{RizviLin2023book}, in
which a window of past inputs and outputs stands in for the state with
no observer; the extended kernel inherits the $2\times2$ block form,
and the complete state-feedback machinery carries over once the
inverse of the Lyapunov matrix is exchanged for the Moore--Penrose
pseudoinverse of its rank-deficient extended analogue---an exchange
shown below to be lossless.

\emph{Statement of novelty.}\ To the best of our knowledge, this is
the first work to (i)~carry out a low-and-high gain design for
discrete-time saturated systems entirely without a model, in state- as
well as output-feedback form; (ii)~establish that one Schur complement
of the learned Q-kernel delivers the Lyapunov matrix, and with it both
a certified, closed-form boost $\rho^{*}$ and the exact stability
limit $\rho_{\max}$, so that saturation avoidance and strict Lyapunov
decrease are guaranteed with no identification, extra learning, or
search; (iii)~establish the \emph{single-substitution
principle} $P^{-1}\!\to\!\bar P^{\dagger}$, under which the entire
state-feedback pipeline transports to output feedback with no loss of
information; and (iv)~identify a structural \emph{cost-rank trade-off}
between gain-boost margin and pole-movement effectiveness, set by the
rank of the cost weighting.

\emph{Relation to the conference version.}\ A preliminary version of
the state-feedback results appeared in the conference
paper~\citep{DonyRizviLin2026CDC}.  The present paper is a substantial
extension: the entire output-feedback framework of
Section~\ref{sec:proposed_of}---the state parameterization, the
extended Q-kernel, and the pseudoinverse theory of
Lemmas~\ref{lemma:Pbar} and~\ref{lemma:Ubar_max} and
Theorem~\ref{thm:of_main}---the single-substitution principle, the
cost-rank/pole-movement analysis of
Section~\ref{sec:pole_movement}, and all output-feedback experiments
are new; the state-feedback material is included in condensed form to
keep the paper self-contained.

\emph{Organization.}\ Section~\ref{sec:problem} reviews the setup and
baseline Q-learning; Sections~\ref{sec:proposed_sf}
and~\ref{sec:proposed_of} develop the state- and output-feedback
extensions; Section~\ref{sec:results} reports results on a fourth-order
benchmark; Section~\ref{sec:conclusion} concludes.

\begin{table*}[t]
\centering
\caption{Comparison with related work.  All methods address actuator
saturation; a check mark denotes the additional property.
``Closed-form $\rho^{*}$'': an analytic high-gain parameter with no
model and no search.}
\label{tab:related_work}
\small
\begin{tabular}{lccccc}
\toprule
Approach & Model-free & Discrete-time & Low-\&-high gain & Output feedback & Closed-form $\rho^{*}$\\
\midrule
Lin \& Saberi (1995)  & $\times$ & $\times$ & $\checkmark$ & $\times$ & $\times$\\
Lin et al.\ (2000)    & $\times$ & $\checkmark$ & $\checkmark$ & $\times$ & $\times$\\
Wang et al.\ (2012)   & $\times$ & $\checkmark$ & $\checkmark$ & $\times$ & $\checkmark$\\
Rizvi \& Lin (2023)   & $\checkmark$ & $\checkmark$ & $\times$ & $\checkmark$ & $\times$\\
\textbf{This work}    & $\checkmark$ & $\checkmark$ & $\checkmark$ & $\checkmark$ & $\checkmark$\\
\bottomrule
\end{tabular}
\end{table*}


%% ======================================================================
\section{Problem Formulation and Background}\label{sec:problem}
%% ======================================================================

\subsection{System Description and Assumptions}

The plant is a discrete-time linear time-invariant system whose input
saturates,
\begin{equation}
x_{k+1} = Ax_k + B\sigma(u_k),
\label{eq:system}
\end{equation}
in which $x_k \in \mathbb{R}^n$ and $u_k \in \mathbb{R}^m$ denote the state and the
input, and $\sigma : \mathbb{R}^m \to \mathbb{R}^m$ clips every input channel at the
level $b > 0$.

\begin{assumption}[ANCBC]\label{assum:ANCBC}
The pair $(A, B)$ is stabilizable and all eigenvalues of $A$ lie on or
inside the unit circle.
\end{assumption}

\subsection{Parameterized DARE and Low Gain Design}

Low-gain feedback is generated by the $\gamma$-parameterized
discrete-time algebraic Riccati equation (DARE),
\begin{equation}
A^\top\! PA - P + \gamma I
  - A^\top\! PB(B^\top\! PB + I)^{-1}B^\top\! PA = 0,
\label{eq:DARE}
\end{equation}
with low-gain parameter $\gamma \in (0,1]$; the associated feedback
gain reads
\begin{equation}
K^*(\gamma)
   = -(B^\top P^*(\gamma) B + I)^{-1} B^\top P^*(\gamma) A.
\label{eq:optimal_gain}
\end{equation}
Assumption~\ref{assum:ANCBC} guarantees a unique solution
$P^*(\gamma) \succ 0$, Schur stability of $A + BK^*(\gamma)$
throughout $(0,1]$, and the defining low-gain
limits~\citep{Lin1998}:
\begin{align}
\lim_{\gamma \to 0} P^*(\gamma) &= 0, \label{eq:P_limit} \\
\lim_{\gamma \to 0} K^*(\gamma) &= 0. \label{eq:K_limit}
\end{align}


\subsection{Q-Function Formulation}

Equation~\eqref{eq:DARE} is the Riccati equation of the LQR problem
whose running cost is
\begin{equation}
r(x_k, u_k) = \gamma \|x_k\|^2 + \|u_k\|^2.
\label{eq:stage_cost}
\end{equation}
The associated Q-function is quadratic,
\begin{equation}
Q(x_k, u_k) = z_k^\top H(\gamma)\, z_k,
\quad z_k = \begin{bmatrix} x_k \\ u_k \end{bmatrix},
\label{eq:Q_function}
\end{equation}
its kernel carrying the block form
\begin{equation}
H(\gamma)
  = \begin{bmatrix} H_{xx} & H_{xu} \\ H_{ux} & H_{uu} \end{bmatrix}
  = \begin{bmatrix}
        \gamma I + A^\top\! PA & A^\top\! PB \\
        B^\top\! PA            & I + B^\top\! PB
    \end{bmatrix}.
\label{eq:H_matrix}
\end{equation}
The optimal gain is read off as
\begin{equation}
K^*(\gamma) = -H_{uu}^{-1} H_{ux},
\label{eq:K_from_H}
\end{equation}
in agreement with~\eqref{eq:optimal_gain}.

\subsection{Model-Free Q-Learning}\label{sec:qlearning}

The kernel obeys the Bellman relation
\begin{equation}
z_k^\top H z_k = \gamma\|x_k\|^2 + \|u_k\|^2 + z_{k+1}^\top H z_{k+1},
\label{eq:Bellman}
\end{equation}
an equation linear in the $s = (n{+}m)(n{+}m{+}1)/2$ distinct entries
of the symmetric $H$, so $L \geq s$ excited tuples
$\{(x_k, u_k, x_{k+1})\}$ determine $H$ by least squares with $(A,B)$
never appearing~\citep{RizviLin2018CDC,RizviLin2023book}; policy
iteration~\citep{Kleinman1968} interleaves this evaluation with the
update $K = -H_{uu}^{-1}H_{ux}$.

Once learning has converged, the baseline design certifies the
actuator constraint only along the realized trajectory,
\begin{equation}
\max_{k = 0,1,\ldots} \|K^*(\gamma) x_k\|_\infty \leq b.
\label{eq:pointwise_check}
\end{equation}
halving $\gamma$ and relearning whenever~\eqref{eq:pointwise_check}
fails.  A contracting trajectory, however, says nothing about the
worst case over the invariant set, so $\gamma$ ends up far smaller
than necessary and the loop correspondingly slow.

\subsection{Notation}\label{sec:notation}

Throughout, bars mark output-feedback quantities
(Section~\ref{sec:proposed_of}): each state-feedback symbol $X$ has
the analogue $\bar X$, and the scalar parameters
$\rho_{\mathrm{sat}}, \rho_{\mathrm{opt}}, \rho^{*}$ are shared, their
defining formulas carrying over verbatim.



%% ===================== END PART 2 (Sections I and II) =================
%% Next: Part 3 -- Section III (State-Feedback Low-and-High Gain
%% Extension): Lemmas, Propositions, Theorem, Algorithm, and figures.

%% =====================================================================
\section{State-Feedback Low-and-High Gain Extension}
\label{sec:proposed_sf}
%% =====================================================================

We now build the high-gain layer from the learned kernel and nothing
else.  Fig.~\ref{fig:pipeline} traces the route: the blocks of $H$
deliver $P$ and $K$; two boost limits are then evaluated
independently---$\rho_{\mathrm{sat}}$ from the actuator geometry and
$\rho_{\mathrm{opt}}$ from the Lyapunov decrease---and the smaller of
the two fixes $K_{\mathrm{high}} = (1+\rho^*)K$.  At no point do
$(A,B)$ enter.

\begin{figure}[!htb]
\centering
\includegraphics[width=\linewidth]{fig/pipeline_figure_v2.pdf}
\caption{The model-free boost pipeline.  Two independent branches
leave the Q-kernel $H$: the upper one passes through $P$, $c$, and
$U_{\max}$ to the saturation limit $\rho_{\mathrm{sat}}$; the lower
one reads $\lambda_{\max}(H_{uu})$ and returns the decrease-margin
maximizer $\rho_{\mathrm{opt}}$---to be distinguished from the
stability limit $\rho_{\max}$ of~\eqref{eq:rho_max}.  The smaller of
the two fixes $\rho^{*}$ and with it
$K_{\mathrm{high}} = (1+\rho^{*})K$.}
\label{fig:pipeline}
\end{figure}

\subsection{Extraction of the Lyapunov Matrix}\label{sec:P_extraction}

\begin{lemma}[Schur Complement Extraction]\label{lemma:P}
Let $H(\gamma)$ be the Q-function kernel with structure
\eqref{eq:H_matrix}.  Then the DARE solution satisfies
\begin{equation}
P(\gamma) = H_{xx} - H_{xu}\, H_{uu}^{-1}\, H_{ux}.
\label{eq:P_from_H}
\end{equation}
\end{lemma}

\begin{proof}
From the block structure~\eqref{eq:H_matrix},
\begin{align}
H_{xx} &= \gamma I + A^\top PA, \label{eq:Hxx_expand} \\
H_{xu} &= A^\top PB,             \label{eq:Hxu_expand} \\
H_{uu} &= I + B^\top PB.         \label{eq:Huu_expand}
\end{align}
Rewriting the DARE~\eqref{eq:DARE} as
$A^\top PA - A^\top PB(I + B^\top PB)^{-1}B^\top PA = P - \gamma I$
and substituting,
\begin{align}
H_{xx} - H_{xu}\, H_{uu}^{-1}\, H_{ux}
  &= \gamma I + A^\top PA \notag \\
  &\quad - A^\top PB(I + B^\top PB)^{-1}B^\top PA \notag \\
  &= \gamma I + (P - \gamma I) = P. \qedhere
\end{align}
\end{proof}

Because $H$ itself is learned without $(A,B)$,
identity~\eqref{eq:P_from_H} hands the low-gain Lyapunov
matrix~\citep{Lin1998} over directly from data.

\subsection{Maximum Control over the Lyapunov Level Set}
\label{sec:Umax}

For the learned $P$ and a given $x_0$, the sublevel set
\begin{equation}
\mathcal{L}_V(c) = \{x \in \mathbb{R}^n : x^\top P x \leq c\},
\quad c = x_0^\top P x_0,
\end{equation}
is forward invariant under the LQR policy.  Sizing the admissible
boost requires the largest control the feedback can request anywhere
on this set.

\begin{lemma}[Worst-Case Control Magnitude]\label{lemma:Umax}
Given $K \in \mathbb{R}^{m \times n}$, $P > 0$, and $c = x_0^\top Px_0$, the
maximum control magnitude over $\mathcal{L}_V(c)$ is
\begin{equation}
U_{\max} = \sqrt{c \cdot \lambda_{\max}(K P^{-1} K^\top)}.
\label{eq:Umax}
\end{equation}
\end{lemma}

\begin{proof}
We seek $U_{\max} = \max_{x^\top Px \leq c} \|Kx\|_2$.

\textit{Step~1 (Coordinate transformation).}  As $P \succ 0$, the
substitution $y = P^{1/2}x$ (so $x = P^{-1/2}y$) turns the ellipsoidal
constraint into a ball:
\begin{equation}
x^\top Px = y^\top P^{-1/2} P P^{-1/2} y = y^\top y = \|y\|^2 \leq c.
\end{equation}

\textit{Step~2 (Quadratic form).}  The objective becomes
\begin{align}
\|Kx\|^2 &= \|KP^{-1/2}y\|^2 \notag \\
          &= y^\top
             \underbrace{P^{-1/2}K^\top K P^{-1/2}}_{\Gamma}\, y.
\end{align}

\textit{Step~3 (Maximization over ball).}  The maximum of
$y^\top \Gamma y$ subject to $\|y\|^2 \leq c$ is
\begin{equation}
\max_{\|y\|^2 \leq c} y^\top \Gamma y = c \cdot \lambda_{\max}(\Gamma),
\end{equation}
the maximizer $y^* = \sqrt{c}\, v_1$ lying along the top eigenvector
$v_1$ of $\Gamma$.

\textit{Step~4 (Eigenvalue invariance).}  Since $MN$ and $NM$ share
their nonzero spectrum for compatibly sized $M$, $N$,
\begin{align}
\lambda_{\max}(\Gamma)
  &= \lambda_{\max}(P^{-1/2}K^\top K P^{-1/2}) \notag \\
  &= \lambda_{\max}(K P^{-1/2} P^{-1/2} K^\top) \notag \\
  &= \lambda_{\max}(K P^{-1} K^\top).
\end{align}
Combining,
$U_{\max} = \sqrt{c \cdot \lambda_{\max}(KP^{-1}K^\top)}$.
\end{proof}

Every ingredient of~\eqref{eq:Umax} is already at hand from $H$ and
$x_0$; the geometry of the proof is sketched in
Fig.~\ref{fig:ellipsoid}.

\begin{figure}[!htb]
\centering
\includegraphics[width=\linewidth]{fig/ellipsoid_transform.pdf}
\caption{Lemma~\ref{lemma:Umax}: $y = P^{1/2}x$ maps the level-set
ellipsoid to the ball $\|y\|^2 \leq c$; the maximizer aligns with the
leading eigenvector of $\Gamma$.}
\label{fig:ellipsoid}
\end{figure}



\subsection{Saturation Bound}\label{sec:rho_sat}

In the spirit of \citet{LinSaberi1995}, the learned low-gain action
$u_{\mathrm{L}} = Kx_k$ is augmented by the high-gain term
$u_{\mathrm{H}} = \rho\,Kx_k$, $\rho \geq 0$:
\begin{equation}
u = u_{\mathrm{L}} + u_{\mathrm{H}}
  = (1+\rho)\, Kx_k = K_{\mathrm{high}}\, x_k,
\label{eq:u_combined}
\end{equation}
so that $K_{\mathrm{high}} = (1+\rho)K$ and $\rho = 0$ returns the
baseline.  Worst-case magnitudes scale with $(1+\rho)$,
\begin{equation}
\max_{x \in \mathcal{L}_V(c)} \|K_{\mathrm{high}}\, x\|_2
  = (1+\rho)\, U_{\max},
\label{eq:boosted_Umax}
\end{equation}
and imposing $(1+\rho)\, U_{\max} \leq b$ gives
\begin{equation}
\rho \leq \frac{b}{U_{\max}} - 1.
\label{eq:sat_derivation}
\end{equation}

\begin{proposition}[Saturation Bound]\label{prop:rho_sat}
Define
\begin{equation}
\rho_{\mathrm{sat}} = \frac{b}{U_{\max}} - 1.
\label{eq:rho_sat}
\end{equation}
If $\rho_{\mathrm{sat}} > 0$ (i.e., $U_{\max} < b$), then for any
$\rho \in [0, \rho_{\mathrm{sat}}]$,
\begin{equation}
\|(1+\rho)Kx\|_2 \leq b, \quad \forall\, x \in \mathcal{L}_V(c).
\end{equation}
\end{proposition}

\begin{proof}
For any $x \in \mathcal{L}_V(c)$, Lemma~\ref{lemma:Umax} gives
$\|Kx\|_2 \leq U_{\max}$.  Then
\begin{align}
\|(1+\rho)Kx\|_2 &= (1+\rho)\|Kx\|_2 \notag \\
                  &\leq (1+\rho)\, U_{\max} \notag \\
                  &\leq \frac{b}{U_{\max}} \cdot U_{\max} = b,
\end{align}
the final step using $\rho \leq \rho_{\mathrm{sat}}$; the bound is
tight at $\rho = \rho_{\mathrm{sat}}$.
\end{proof}

\subsection{Lyapunov Decrease Bound}\label{sec:lyapunov}

Respecting the actuator limit is not the same as remaining stable:
too aggressive a boost can drive $V(x) = x^\top Px$ upward.  We
therefore characterize, for arbitrary $m \geq 1$ and purely in terms
of the kernel, the range of $\rho$ on which $V$ strictly decreases
under $u_k = (1+\rho)Kx_k$ of~\eqref{eq:u_combined}.

The two Bellman facts $V(x_k) = Q(x_k,Kx_k)$ and
$V(x_{k+1}) = Q(x_k,u_k) - \gamma\|x_k\|^2 - \|u_k\|^2$, differenced
at $u_k = (1+\rho)Kx_k$, give
\begin{equation}
V(x_k) - V(x_{k+1})
   = \Delta Q + \gamma\|x_k\|^2 + (1+\rho)^2\|Kx_k\|^2,
\label{eq:V_diff}
\end{equation}
in which, expanding the quadratic form of $Q$ and invoking
$H_{xu} = -K^\top H_{uu}$ (a consequence of $K = -H_{uu}^{-1}H_{ux}$
together with symmetry),
\begin{align}
\Delta Q
   &= Q(x,Kx) - Q\bigl(x,(1{+}\rho)Kx\bigr) \notag\\
   &= -2\rho\,x^\top\! H_{xu}Kx \notag\\
   &\quad {}+ \bigl[1-(1{+}\rho)^2\bigr](Kx)^\top\! H_{uu}(Kx) \notag\\
   &= -\rho^2\,(Kx)^\top\! H_{uu}\,(Kx).
\label{eq:DeltaQ_final}
\end{align}
Insertion into~\eqref{eq:V_diff}, with the $\|Kx_k\|^2$ term
absorbed, gives
\begin{equation}
V(x_k) - V(x_{k+1})
   = \gamma\|x_k\|^2 + (Kx_k)^\top \Lambda(\rho)\,(Kx_k),
\label{eq:lyap_decrease}
\end{equation}
with the \emph{Lyapunov-decrease indicator}
\begin{equation}
\boxed{\,\Lambda(\rho) \triangleq (1+\rho)^2 I_m - \rho^2 H_{uu}.\,}
\label{eq:Lambda}
\end{equation}
Positivity of $\Lambda(\rho)$ is the stability requirement.  Every
eigenvalue $\lambda_i(H_{uu})$ contributes
$(1+\rho)^2 - \rho^2\lambda_i(H_{uu})$ to the spectrum of
$\Lambda(\rho)$, and since $H_{uu} = I + B^\top PB$ forces
$\bar\lambda \triangleq \lambda_{\max}(H_{uu}) > 1$, the smallest of
these is
\begin{equation}
\lambda_{\min}\!\big(\Lambda(\rho)\big)
   = (1+\rho)^2 - \rho^2\bar\lambda \;\triangleq\; f(\rho),
\label{eq:Lambda_min}
\end{equation}
a concave parabola through $f(0) = 1$ (Fig.~\ref{fig:lambda}).  The
parabola carries two separate pieces of information, and the
distinction matters.  Its positive root, obtained from
$1 + \rho = \rho\sqrt{\bar\lambda}$, marks where the certificate
expires:
\begin{equation}
\Lambda(\rho) \succ 0
   \;\iff\;
   \rho < \rho_{\max} \triangleq \frac{1}{\sqrt{\bar\lambda} - 1};
\label{eq:rho_max}
\end{equation}
only there---not at $\rho_{\mathrm{opt}}$ below---is stability at
stake.  Its vertex, from $df/d\rho = 0$, locates the strongest
certificate and its height,
\begin{equation}
\rho_{\mathrm{opt}} \triangleq \frac{1}{\bar\lambda - 1},
\qquad
f(\rho_{\mathrm{opt}}) = \frac{\bar\lambda}{\bar\lambda - 1}
   = 1 + \rho_{\mathrm{opt}}.
\label{eq:rho_opt}
\end{equation}
As $\bar\lambda - 1 > \sqrt{\bar\lambda} - 1$,
\begin{equation}
0 < \rho_{\mathrm{opt}} < \rho_{\max},
\label{eq:rho_order}
\end{equation}
the maximizer is interior to the admissible range, and monotonicity
of $f$ on $[0,\rho_{\mathrm{opt}}]$ yields $f(\rho) \geq f(0) = 1$
there.  Selecting $\rho_{\mathrm{opt}}$ is thus a \emph{performance}
decision---it delivers the strongest certified decrease---rather than
a stability necessity: Case~1 of Section~\ref{sec:results} has
$\rho_{\mathrm{opt}} = 2.31$ versus $\rho_{\max} = 5.07$, and every
intermediate $\rho$ remains stable, with a certificate thinning to
zero toward $\rho_{\max}$.

\begin{figure}[!htb]
\centering
\includegraphics[width=\linewidth]{fig/lambda_figure_fixed.pdf}
\caption{Spectrum of $\Lambda(\rho)$: every $\lambda_i(H_{uu})$
traces the parabola $(1+\rho)^2 - \rho^2\lambda_i(H_{uu})$, all
passing through $1$ at $\rho = 0$, with
$\bar\lambda = \lambda_{\max}(H_{uu})$ the binding one.  The
certificate holds precisely on $[0,\rho_{\max})$,
$\rho_{\max} = 1/(\sqrt{\bar\lambda}-1)$~\eqref{eq:rho_max}, while the
design operates at the interior vertex
$\rho_{\mathrm{opt}} = 1/(\bar\lambda-1)$~\eqref{eq:rho_opt}; for
$m = 1$ a single parabola survives.}
\label{fig:lambda}
\end{figure}

\begin{proposition}[Selection of the High Gain Parameter]\label{prop:rho_star}
Define
\begin{equation}
\rho^* = \max\!\big(0,\; \min(\rho_{\mathrm{sat}},\,
        \rho_{\mathrm{opt}})\big),
\quad \rho_{\mathrm{opt}} = \frac{1}{\lambda_{\max}(H_{uu})-1}.
\label{eq:rho_star}
\end{equation}
Then $K_{\mathrm{high}} = (1+\rho^*)K$ satisfies:
\begin{enumerate}
\item[(a)] Saturation avoidance:
   $\|K_{\mathrm{high}} x\| \leq b$, $\forall\, x \in \mathcal{L}_V(c)$;
\item[(b)] Lyapunov decrease: $\Lambda(\rho^*) \succ 0$, giving
   $V(x_k) - V(x_{k+1})
     = \gamma\|x_k\|^2
     + (Kx_k)^\top \Lambda(\rho^*)(Kx_k) > 0$
   for all $x_k \neq 0$;
\item[(c)] Model-free: all quantities computed from $H$ and $x_0$.
\end{enumerate}
\end{proposition}

\subsection{Stability Guarantee}\label{sec:stability}

\begin{theorem}[Model-Free High-Gain Stabilization]\label{thm:main}
Under Assumption~\ref{assum:ANCBC}, let $H(\gamma)$ be the converged
Q-kernel, with $K = -H_{uu}^{-1} H_{ux}$,
$P = H_{xx} - H_{xu} H_{uu}^{-1} H_{ux}$, $c = x_0^\top P x_0$,
$U_{\max} = \sqrt{c\,\lambda_{\max}(KP^{-1}K^\top)}$,
$\rho_{\mathrm{sat}} = b/U_{\max} - 1$,
$\rho_{\mathrm{opt}} = 1/(\lambda_{\max}(H_{uu}) - 1)$, and
$\rho^{*} = \max\!\big(0, \min(\rho_{\mathrm{sat}},
\rho_{\mathrm{opt}})\big)$ (equivalently, Algorithm~\ref{alg:sf}).  If
$\rho^{*} > 0$, the controller $K_{\mathrm{high}} = (1+\rho^{*})K$
guarantees:
\begin{enumerate}
\item[(a)] \textbf{Saturation avoidance:}
   $\|K_{\mathrm{high}} x_k\|_\infty \leq b$ for all
   $x_k \in \mathcal{L}_V(c)$;
\item[(b)] \textbf{Asymptotic stability:}
   $V(x_{k+1}) < V(x_k)$ for all $x_k \neq 0$ in $\mathcal{L}_V(c)$;
\item[(c)] \textbf{Positive invariance:} $\mathcal{L}_V(c)$ is
   positively invariant.
\end{enumerate}
Moreover, for any bounded set of initial conditions, there exists a
sufficiently small $\gamma$ with $\rho^{*} > 0$ (semi-global
stabilization).
\end{theorem}

\begin{proof}
\emph{(a) Saturation avoidance.}  Fix $x_k \in \mathcal{L}_V(c)$.
Lemma~\ref{lemma:Umax} bounds $\|Kx_k\|_2$ by $U_{\max}$, and with
$u_k = (1+\rho^*)Kx_k$ and $\rho^* \leq \rho_{\mathrm{sat}} =
b/U_{\max}-1$,
\begin{equation}
\|u_k\|_\infty \leq \|u_k\|_2 = (1+\rho^*)\|Kx_k\|_2
   \leq (1+\rho_{\mathrm{sat}})U_{\max} = b,
\end{equation}
and the saturation never engages on $\mathcal{L}_V(c)$.

\emph{(b) Asymptotic stability.}  Saturation being inactive, the
identity of Section~\ref{sec:lyapunov} applies:
$V(x_k) - V(x_{k+1}) =
\gamma\|x_k\|^2 + (Kx_k)^\top\Lambda(\rho^*)(Kx_k)$.  The
ordering~\eqref{eq:rho_order} places $\rho^* \leq \rho_{\mathrm{opt}}$
inside the monotone range of $f$, so by~\eqref{eq:Lambda_min}
$\lambda_{\min}(\Lambda(\rho^*)) = f(\rho^*) \geq f(0) = 1$ and
$\Lambda(\rho^*) \succeq I_m \succ 0$; as $\gamma > 0$, the right-hand
side is positive whenever $x_k \neq 0$, hence
$V(x_{k+1}) < V(x_k)$.

\emph{(c) Positive invariance.}  Monotonicity of $V$ gives
$V(x_{k+1}) < V(x_k) \leq c$, hence $x_{k+1} \in \mathcal{L}_V(c)$.

\emph{Semi-global property.}  The limit $P \to 0$ does not by itself
force $U_{\max} \to 0$, because~\eqref{eq:Umax} contains $P^{-1}$ as
well; we bound $U_{\max}$ through $c$ instead.  At the optimal gain,
with $A_{\mathrm{cl}} = A + BK$, the DARE~\eqref{eq:DARE} takes the
closed-loop form
\begin{equation}
A_{\mathrm{cl}}^\top P A_{\mathrm{cl}} - P
   = -\bigl(\gamma I + K^\top\! K\bigr),
\label{eq:cl_lyap}
\end{equation}
whence $K^\top\! K \preceq P$ and
$\|Kx\|_2^2 \leq x^\top\! Px \leq c$ across the level set, so that
$U_{\max} \leq \sqrt{c} = \sqrt{x_0^\top P(\gamma)\,x_0}$ and
$\lambda_{\max}(KP^{-1}K^\top) \leq 1$.  As $c \to 0$
by~\eqref{eq:P_limit}, $U_{\max}(\gamma)$ vanishes uniformly on
bounded sets of initial states, $\rho_{\mathrm{sat}} \to \infty$, and
$\rho^* = \rho_{\mathrm{opt}} > 0$ is eventually reached---located
automatically by the $\gamma$-schedule.
\end{proof}

\subsection{Algorithm}\label{sec:algorithm}

Algorithm~\ref{alg:sf} collects the procedure; Lines~4--11 are
algebra on the converged kernel, so the boost costs no further data or
learning.  The $\gamma$-schedule accepts only when
$\rho_{\mathrm{sat}} > 0$, i.e., $U_{\max} < b$---a condition over the
whole invariant set that, by
$\max_k\|Kx_k\|_\infty \leq U_{\max}$,
subsumes the trajectory test~\eqref{eq:pointwise_check} without the
converse holding.

\begin{algorithm}[!htb]
\caption{Model-Free Low-and-High Gain Q-Learning (State Feedback)}
\label{alg:sf}
{\footnotesize
\begin{algorithmic}[1]
\Require $x_0$, limit $b$, initial $\gamma$, tolerance $\epsilon$
\Ensure  high-gain controller $K_{\mathrm{high}}$
\State Initialize stabilizing $K^0$;\ $\gamma_c \leftarrow \gamma$
\Repeat
    \State Run PI at $\gamma_c$ until $\|H^j - H^{j-1}\| < \epsilon$
    \State $K \leftarrow -H_{uu}^{-1}H_{ux}$
    \State $P \leftarrow H_{xx} - H_{xu}H_{uu}^{-1}H_{ux}$
    \State $c \leftarrow x_0^\top Px_0$
    \State $U_{\max} \leftarrow \sqrt{c\,\lambda_{\max}(KP^{-1}K^\top)}$
    \State $\rho_{\mathrm{sat}} \leftarrow b/U_{\max} - 1$
    \State $\rho_{\mathrm{opt}} \leftarrow 1/(\lambda_{\max}(H_{uu})-1)$
    \State $\rho^* \leftarrow \max(0,\min(\rho_{\mathrm{sat}},\rho_{\mathrm{opt}}))$
    \State $K_{\mathrm{high}} \leftarrow (1+\rho^*)K$
    \If{$\rho_{\mathrm{sat}} \leq 0$}\Comment{$U_{\max}\!\geq\! b$: invariant-set test}
    \State $\gamma_c \leftarrow \gamma_c/2$ \EndIf
\Until{$\rho_{\mathrm{sat}} > 0$}
\State \Return $K_{\mathrm{high}}$, $\rho^*$, $\gamma_c$
\end{algorithmic}}
\end{algorithm}

%% =====================================================================
\section{Output-Feedback Low-and-High Gain Extension}\label{sec:proposed_of}
%% =====================================================================

We extend the framework of Section~\ref{sec:proposed_sf} to output
feedback, where only the input--output history is available and $x_k$
is not measured.  The entire SF pipeline transports once $x_k$ is
replaced by an extended history vector and $P^{-1}$ by the
Moore--Penrose pseudoinverse $\bar P^{\dagger}$, with $(A,B,C)$ never
appearing.  Two lemmas carry the new content: a Schur identity
extracting the rank-deficient extended Lyapunov matrix $\bar P$ from the
learned kernel (Lemma~\ref{lemma:Pbar}), and a level-set identity
showing the worst-case-control formula survives the rank deficiency
(Lemma~\ref{lemma:Ubar_max}).  We use the
state-parameterization of~\citet{RizviLin2023book}, which requires
neither a separate observer nor discounting.

\subsection{System Description with Output Measurements}
\label{sec:of_system}

We augment the system~\eqref{eq:system} with a measurement equation,
\begin{align}
x_{k+1} &= A\,x_k + B\,\sigma(u_k), \label{eq:of_system_x}\\
y_k     &= C\,x_k, \label{eq:of_system_y}
\end{align}
where $y_k \in \mathbb{R}^p$ is the measured output and the saturation function
$\sigma$ acts component-wise with limit $b > 0$ as before.

\begin{assumption}[Observability]\label{assum:obs}
The pair $(A, C)$ is observable.
\end{assumption}

Together with Assumption~\ref{assum:ANCBC}, this is the standard
setting for output-feedback Q-learning of saturated systems.

\subsection{State Parameterization via Input/Output Histories}
\label{sec:of_param}

Without the system matrices a Luenberger observer is unavailable, but
stacking past inputs and outputs with a dead-beat parameterization
recovers the state exactly~\citep{RizviLin2023book}; the result we need
is the following.

\begin{lemma}[State parameterization]
\label{lem:state_param}
Let $(A, C)$ be observable and let $N \leq n$ be an upper bound on the
observability index of $(A, C)$.  Define the past-input and
past-output histories
\begin{align}
\bar u_{k-1, k-N}
   &= \begin{bmatrix} u_{k-1}^{\!\top} & u_{k-2}^{\!\top}
                      & \cdots & u_{k-N}^{\!\top} \end{bmatrix}^{\!\top}
      \in \mathbb{R}^{mN}, \label{eq:past_u}\\[2pt]
\bar y_{k-1, k-N}
   &= \begin{bmatrix} y_{k-1}^{\!\top} & y_{k-2}^{\!\top}
                      & \cdots & y_{k-N}^{\!\top} \end{bmatrix}^{\!\top}
      \in \mathbb{R}^{pN}. \label{eq:past_y}
\end{align}
Then there exist constant matrices $W_u \in \mathbb{R}^{n \times mN}$ and
$W_y \in \mathbb{R}^{n \times pN}$ such that, for all $k \geq N$,
\begin{equation}
x_k = W_u\,\bar u_{k-1, k-N} + W_y\,\bar y_{k-1, k-N}.
\label{eq:state_param}
\end{equation}
The matrices $W_u, W_y$ depend on $(A, B, C)$ and on the dead-beat
horizon $N$, but the relation~\eqref{eq:state_param} is exact (not
asymptotic) for $k \geq N$.
\end{lemma}

Combine the histories into an \emph{extended state}
\begin{equation}
\xi_k \;\triangleq\;
   \begin{bmatrix} \bar u_{k-1, k-N} \\[2pt] \bar y_{k-1, k-N} \end{bmatrix}
   \in \mathbb{R}^{(m+p)N},
\label{eq:xi_def}
\end{equation}
and the matrices into a \emph{state-recovery matrix}
\begin{equation}
W \;\triangleq\;
   \begin{bmatrix} W_u & W_y \end{bmatrix}
   \in \mathbb{R}^{n \times (m+p)N}.
\label{eq:W_def}
\end{equation}
so Lemma~\ref{lem:state_param} reads
\begin{equation}
x_k = W\,\xi_k, \quad k \geq N.
\label{eq:state_param_compact}
\end{equation}
Since $(m+p)N \gg n$ in general, $W$ is wide with
$\mathrm{rank}\,W \leq n$; under the rank conditions
of~{\citep[Theorems~2.3, 2.4]{RizviLin2023book}} it has full row rank
$n$, which we assume.

\begin{assumption}[Full row rank of $W$]\label{assum:W_rank}
The state-recovery matrix $W \in \mathbb{R}^{n \times (m+p)N}$ has full row
rank, $\mathrm{rank}\,W = n$.
\end{assumption}

\subsection{Output-Feedback Parameterized DARE and Q-Function}
\label{sec:of_q}

The conventional state-feedback LQR cost~\eqref{eq:stage_cost} weights
the full state, which is unavailable.  Following~\citep{RizviLin2023book},
we instead weight only the measured output:
\begin{equation}
r(x_k, u_k)
   = \gamma\,y_k^{\!\top}\!y_k + u_k^{\!\top}\!u_k
   = \gamma\,x_k^{\!\top} C^{\!\top}\!C\,x_k + u_k^{\!\top}\!u_k.
\label{eq:of_stage_cost}
\end{equation}
The corresponding state weighting matrix is
$Q = \gamma\,C^{\!\top}\!C \in \mathbb{R}^{n \times n}$, which has rank at most
$p$ (and rank exactly one for the single-output case $p = 1$).  The OF
parameterized DARE is then
\begin{equation}
\begin{aligned}
&A^{\!\top}\!P A - P + \gamma\,C^{\!\top}\!C \\
&\qquad - A^{\!\top}\!PB\,(B^{\!\top}\!PB + I)^{-1}\,B^{\!\top}\!PA = 0,
\end{aligned}
\label{eq:of_dare}
\end{equation}
which has the same structure as the SF DARE~\eqref{eq:DARE} with
$\gamma I$ replaced by $\gamma\,C^{\!\top}\!C$.  Under
Assumptions~\ref{assum:ANCBC}, \ref{assum:obs}, and~\ref{assum:W_rank},
the unique positive definite solution $P^{*}(\gamma)$ exists, the
closed-loop matrix $A + B K^{*}(\gamma)$ is Schur stable for all
$\gamma \in (0, 1]$, and the low-gain limit~\eqref{eq:P_limit}--%
\eqref{eq:K_limit} continues to hold.

\paragraph{Extended Q-function.}
Substituting $x_k = W\,\xi_k$
\eqref{eq:state_param_compact} into the SF Q-function
$Q = z_k^{\!\top}\!H z_k$ \eqref{eq:Q_function}, i.e., $z_k = T\bar z_k$
with $T = \operatorname{diag}(W, I_m)$ and
$\bar z_k = [\xi_k^{\!\top}, u_k^{\!\top}]^{\!\top}$, gives for
$k \geq N$ the extended form $Q = \bar z_k^{\!\top}\bar H\bar z_k$ with
\begin{equation}
\bar H \triangleq T^{\!\top} H\, T
   = \begin{bmatrix} \bar H_{\xi\xi} & \bar H_{\xi u} \\
                     \bar H_{u\xi}   & \bar H_{uu} \end{bmatrix},
\label{eq:Hbar_def}
\end{equation}
the \emph{extended Q-kernel}, whose blocks are the SF blocks with each
$x$-factor transformed by $W$ (the pure-$u$ block unchanged):
\begin{align}
\bar H_{\xi\xi}
   &= W^{\!\top}\!\bigl(\gamma\,C^{\!\top}\!C + A^{\!\top}\!PA\bigr)W,
   \label{eq:Hbar_xixi}\\
\bar H_{\xi u}
   &= W^{\!\top}\!A^{\!\top}\!PB = \bar H_{u\xi}^{\!\top},\\
\bar H_{uu}
   &= I + B^{\!\top}\!PB. \label{eq:Hbar_uu}
\end{align}

Three features are immediate from~\eqref{eq:Hbar_xixi}--\eqref{eq:Hbar_uu}.
First, $\bar H$ has the \emph{same $2\times2$ block form} as the SF
kernel~\eqref{eq:H_matrix} with $\xi_k$ in place of $x_k$, so the
substitution $x\mapsto\xi$, $H\mapsto\bar H$ carries the SF identities to
the OF setting.  Second, $\bar H_{uu} = I + B^{\!\top}\!PB$ is identical
in \emph{form} to the SF $H_{uu}$---which is what lets the entire
Lyapunov-decrease analysis transport unchanged---but not in \emph{value}:
here $P$ solves the OF DARE~\eqref{eq:of_dare} with rank-deficient cost
$Q=\gamma\,C^{\!\top}\!C$, and Riccati monotonicity
($\gamma I \succeq \gamma\,C^{\!\top}\!C \Rightarrow P^{SF}\succeq P^{OF}$)
predicts $\rho_{\mathrm{opt}}^{OF} \geq \rho_{\mathrm{opt}}^{SF}$,
confirmed in Section~\ref{sec:results}.  Third,
$\bar H_{\xi\xi} = W^{\!\top}(\gamma\,C^{\!\top}\!C + A^{\!\top}\!PA)W$ is
rank-deficient---rank at most $n$ in dimension $(m+p)N$, as it factors
through $W$---the one technical complication the section addresses.

The optimal OF gain, obtained from $\partial \bar Q / \partial u = 0$,
is
\begin{equation}
\bar K^{*} \;=\; -\bar H_{uu}^{-1}\,\bar H_{u\xi}
   \;\in\; \mathbb{R}^{m \times (m+p)N},
\label{eq:Kbar_opt}
\end{equation}
which is the OF counterpart of the SF formula~\eqref{eq:K_from_H}.
The control law is $u_k = \bar K^{*}\xi_k$, expressed entirely in
terms of the input/output history---the state $x_k$ does not appear.

\paragraph{Model-free learning of $\bar H$.}
The OF Q-function satisfies the Bellman equation
\begin{equation}
\bar z_k^{\!\top}\,\bar H\,\bar z_k
  = \gamma\,y_k^{\!\top}\!y_k + u_k^{\!\top}\!u_k
    + \bar z_{k+1}^{\!\top}\,\bar H\,\bar z_{k+1},
\quad k \geq N,
\label{eq:of_bellman}
\end{equation}
which is linear in the unique entries of $\bar H$.  Collecting at
least $L \geq \bar s$ data tuples
$\{(y_k, u_k, \xi_k, \xi_{k+1})\}$ with
$\bar s = ((m+p)N + m)\bigl((m+p)N + m + 1\bigr)/2$ and a persistently
exciting probing signal yields $\bar H$ via least squares, exactly as
in the SF case.  The PI/VI algorithms, convergence proofs, and the bias-immunity
result guaranteeing that the learned $\bar H$
matches~\eqref{eq:Hbar_def} \emph{exactly} are given
in~\citep[Ch.~2]{RizviLin2023book}; we treat the learned kernel as
exact in what follows.

The full high-gain pipeline (Fig.~\ref{fig:of_pipeline}) mirrors the SF
case---extract $\bar P$ and $\bar K$ from $\bar H$, take
$\rho^{*} = \max(0,\min(\rho_{\mathrm{sat}}, \rho_{\mathrm{opt}}))$, and
form $\bar K_{\mathrm{high}} = (1+\rho^{*})\bar K$---using only $\bar H$
and $\xi_{\mathrm{init}}$.

\begin{figure}[!htb]
\centering
\includegraphics[width=\linewidth]{fig/others.pdf}
\caption{Output-feedback high-gain pipeline.  The extended kernel
$\bar H$ yields $\bar P$ (Schur complement, Lemma~\ref{lemma:Pbar}),
$\bar K$, and $\rho_{\mathrm{opt}}$; with
$\bar c = \xi_{\mathrm{init}}^{\!\top}\bar P\,\xi_{\mathrm{init}}$ these
give $\rho^{*}$ and $\bar K_{\mathrm{high}} = (1+\rho^{*})\bar K$.  Only
$\bar H$ and $\xi_{\mathrm{init}}$ enter.}
\label{fig:of_pipeline}
\end{figure}

\subsection{Schur Complement Extraction of the Extended Lyapunov Matrix}
\label{sec:of_schur}

\begin{lemma}[OF Schur complement extraction]\label{lemma:Pbar}
Let $\bar H(\gamma)$ be the OF extended Q-kernel with the
structure~\eqref{eq:Hbar_def}--\eqref{eq:Hbar_uu}.  Define the
extended Lyapunov matrix
\begin{equation}
\bar P \;\triangleq\; W^{\!\top}\!P\,W
   \;\in\; \mathbb{R}^{(m+p)N \times (m+p)N},
\label{eq:Pbar_def}
\end{equation}
where $P = P^{*}(\gamma)$ is the OF DARE solution
from~\eqref{eq:of_dare}.  Then $\bar P$ is recoverable from $\bar H$
as the Schur complement of $\bar H_{uu}$ in $\bar H$:
\begin{equation}
\bar P
   = \bar H_{\xi\xi}
     - \bar H_{\xi u}\,\bar H_{uu}^{-1}\,\bar H_{u\xi}.
\label{eq:Pbar_schur}
\end{equation}
Moreover, $\bar P$ is positive semi-definite with
$\mathrm{rank}\,\bar P = n$, and its kernel coincides with the null
space of $W$.
\end{lemma}

\begin{proof}
Substituting~\eqref{eq:Hbar_xixi}--\eqref{eq:Hbar_uu} into
\eqref{eq:Pbar_schur} and applying the OF DARE~\eqref{eq:of_dare},
\begin{align}
&\bar H_{\xi\xi} - \bar H_{\xi u}\bar H_{uu}^{-1}\bar H_{u\xi} \notag\\
&\; = W^{\!\top}\!\bigl[\gamma C^{\!\top}\!C + A^{\!\top}\!PA \notag\\
&\qquad\; - A^{\!\top}\!PB(I{+}B^{\!\top}\!PB)^{-1}B^{\!\top}\!PA\bigr]W \notag\\
&\; = W^{\!\top}\!\bigl[\gamma C^{\!\top}\!C + (P - \gamma C^{\!\top}\!C)\bigr]W \notag\\
&\; = W^{\!\top}\!P\,W = \bar P.
\end{align}
Since $P \succ 0$ and $\mathrm{rank}\,W = n$
(Assumptions~\ref{assum:ANCBC}, \ref{assum:W_rank}),
$\bar P = W^{\!\top}\!PW \succeq 0$ with $\mathrm{rank}\,\bar P = n$ and
$\ker\bar P = \ker W$.
\end{proof}

\subsection{Worst-Case Control over the Extended Lyapunov Level Set}
\label{sec:of_umax}

The extended Lyapunov level set is
\begin{equation}
\bar{\mathcal L}_V(\bar c)
   \;\triangleq\;
   \bigl\{\xi \in \mathbb{R}^{(m+p)N} :
          \xi^{\!\top}\!\bar P\,\xi \leq \bar c \bigr\},
\quad
\bar c = \xi_{\mathrm{init}}^{\!\top}\,\bar P\,\xi_{\mathrm{init}},
\label{eq:Lbar_def}
\end{equation}
with $\xi_{\mathrm{init}} = \xi_N$ the extended state at the start of
the post-learning horizon.  Since $\bar P$ has rank $n$ in dimension
$(m+p)N$, this set is not an ellipsoid but an unbounded
\emph{cylinder}: an ellipsoid in the range of $W^{\!\top}$ extruded
freely along $\ker W$ (Fig.~\ref{fig:of_ellipsoid}, left).  The
worst-case control is nonetheless finite because $\bar K$ annihilates
the unbounded directions: with $\bar K = -\bar H_{uu}^{-1}\bar H_{u\xi}$
and $\bar H_{u\xi} = B^{\!\top}\!PAW$,
\begin{equation}
\bar K\,\xi
   = -\bar H_{uu}^{-1}\,B^{\!\top}\!P A\,W\,\xi
   = K\,(W\,\xi) = K\,x
   \quad (x = W\,\xi),
\label{eq:Kbar_collapse}
\end{equation}
so $\bar K\,\xi = 0$ whenever $W\,\xi = 0$, and the supremum is attained
on the bounded cross section.

\begin{lemma}[OF worst-case control magnitude]\label{lemma:Ubar_max}
Let $\bar K \in \mathbb{R}^{m \times (m+p)N}$, $\bar P \succeq 0$ with
$\mathrm{rank}\,\bar P = n$, and
$\bar c = \xi_{\mathrm{init}}^{\!\top}\,\bar P\,\xi_{\mathrm{init}}$.
Then the supremum of $\|\bar K\,\xi\|_2$ over the extended level
set~\eqref{eq:Lbar_def} is finite and equals
\begin{equation}
\bar U_{\max}
   = \sqrt{\,\bar c\cdot
       \lambda_{\max}\!\bigl(\bar K\,\bar P^{\dagger}\,\bar K^{\!\top}\bigr)\,},
\label{eq:Ubar_max}
\end{equation}
where $\bar P^{\dagger}$ is the Moore--Penrose pseudoinverse of
$\bar P$.  Moreover, the OF formula~\eqref{eq:Ubar_max} matches the SF
formula~\eqref{eq:Umax} exactly:
\begin{equation}
\lambda_{\max}\!\bigl(\bar K\,\bar P^{\dagger}\,\bar K^{\!\top}\bigr)
   = \lambda_{\max}\!\bigl(K\,P^{-1}\,K^{\!\top}\bigr),
\label{eq:Ubar_match_sf}
\end{equation}
so the substitution $P^{-1} \to \bar P^{\dagger}$ produces no
information loss.
\end{lemma}

\begin{proof}
\textit{Step~1 (SVD and decomposition).}  Let
$W = U\,\Sigma\,V^{\!\top}$ be the compact SVD, with
$U \in \mathbb{R}^{n\times n}$ orthogonal,
$\Sigma \in \mathbb{R}^{n\times n}$ diagonal and invertible ($W$ has
full row rank $n$), and $V \in \mathbb{R}^{(m+p)N\times n}$ with
orthonormal columns spanning $\mathrm{range}\,W^{\!\top}$; let
$V_{\perp}$ complete an orthonormal basis of $\ker W$.  Decompose
\begin{equation}
\xi = V\,\alpha + V_{\perp}\,\beta,
\quad \alpha = V^{\!\top}\!\xi,
\quad \beta = V_{\perp}^{\!\top}\!\xi.
\label{eq:xi_decomp}
\end{equation}
Then $W\xi = U\Sigma\,\alpha$ and, by~\eqref{eq:Kbar_collapse},
$\bar K\xi = K\,U\Sigma\,\alpha$---both independent of $\beta$---while
\begin{equation}
\xi^{\!\top}\!\bar P\,\xi
   = (W\xi)^{\!\top}\!P\,(W\xi)
   = \alpha^{\!\top}\!S\,\alpha,
\quad S \triangleq \Sigma\,U^{\!\top}\!P\,U\,\Sigma \succ 0.
\label{eq:constraint_alpha}
\end{equation}

\textit{Step~2 (Reduction to the cross section).}  The constraint
depends only on $\alpha$ (a proper ellipsoid, the cylinder cross
section), so the unbounded $\beta$ directions drop out:
\begin{equation}
\sup_{\xi \in \bar{\mathcal L}_V(\bar c)} \|\bar K\,\xi\|_2^2
   = \max_{\alpha^{\!\top}\!S\,\alpha \leq \bar c}
     \|K\,U\,\Sigma\,\alpha\|_2^2.
\label{eq:reduced_problem}
\end{equation}

\textit{Step~3 (Whitening).}  With $\eta = S^{1/2}\alpha$, the
constraint becomes the ball $\|\eta\|^2 \leq \bar c$ and
\begin{equation}
\|K\,U\,\Sigma\,\alpha\|_2^2 = \eta^{\!\top}\Gamma\,\eta,
\quad
\Gamma \triangleq
S^{-1/2}\Sigma\,U^{\!\top}\!K^{\!\top}\!K\,U\,\Sigma\,S^{-1/2},
\label{eq:whitening}
\end{equation}
whose maximum over the ball is $\bar c\,\lambda_{\max}(\Gamma)$; the
three-stage geometry (cylinder $\to$ ellipsoid $\to$ ball) is shown in
Fig.~\ref{fig:of_ellipsoid}.

\textit{Step~4 (Pseudoinverse identity).}  We claim
\begin{equation}
\bar P^{\dagger}
   = V\,\Sigma^{-1}\,U^{\!\top}\!P^{-1}\,U\,\Sigma^{-1}\,V^{\!\top}.
\label{eq:Pbar_dagger}
\end{equation}
Using $\bar P = V\Sigma\,U^{\!\top}\!P\,U\Sigma\,V^{\!\top}$,
$V^{\!\top}\!V = I_n$, and $U^{\!\top}\!U = U U^{\!\top} = I_n$,
\begin{equation}
\bar P\,\bar P^{\dagger}
  = V\Sigma\,U^{\!\top}\!P\,U\Sigma\,
     \Sigma^{-1}U^{\!\top}\!P^{-1}U\,\Sigma^{-1}V^{\!\top}
  = V\,V^{\!\top},
\label{eq:PPdagger}
\end{equation}
the symmetric idempotent projector onto $\mathrm{range}\,V$, from which
the four Moore--Penrose conditions follow (e.g.\
$\bar P\,\bar P^{\dagger}\bar P = V V^{\!\top}\bar P = \bar P$).  The
same cancellations, with $\bar K = K\,U\Sigma\,V^{\!\top}$, give
\begin{align}
\bar K\,\bar P^{\dagger}\,\bar K^{\!\top}
  &= K\,U\Sigma\,(V^{\!\top}\!V)\,\Sigma^{-1}U^{\!\top}\!P^{-1}
     U\,\Sigma^{-1}(V^{\!\top}\!V)\,\Sigma\,U^{\!\top}\!K^{\!\top}
     \notag\\
  &= K\,(U U^{\!\top})\,P^{-1}\,(U U^{\!\top})\,K^{\!\top}
   = K\,P^{-1}\,K^{\!\top}.
\label{eq:KPdaggerK_eq_KPinvK}
\end{align}
Finally, $\Gamma$ and $\bar K\,\bar P^{\dagger}\bar K^{\!\top}$ share
the same nonzero spectrum (the nonzero eigenvalues of $MN$ and $NM$
coincide, applied cyclically), so
$\lambda_{\max}(\Gamma) = \lambda_{\max}(K\,P^{-1}K^{\!\top})$;
with~\eqref{eq:reduced_problem} this gives~\eqref{eq:Ubar_max} and the
state-feedback match~\eqref{eq:Ubar_match_sf}.
\end{proof}

\begin{remark}[The single-substitution principle]\label{rem:substitution}
Lemma~\ref{lemma:Ubar_max} is the analytical centerpiece: the OF
worst-case-control formula equals the SF formula under the single,
information-preserving substitution $P^{-1}\to\bar P^{\dagger}$; every
downstream quantity is its SF counterpart under the same substitution,
so Sections~\ref{sec:of_sat}--\ref{sec:of_stability} are brief.
\end{remark}

\begin{figure}[!htb]
\centering
\includegraphics[width=\linewidth]{figof/of_ellipsoid_transform (1).pdf}
\caption{Lemma~\ref{lemma:Ubar_max}: the level set is a
rank-deficient cylinder---an ellipsoid in $\mathrm{range}\,W^{\!\top}$
extruded along $\ker W$, where $\bar K$ vanishes (left); projection
onto $\alpha$ collapses it to a proper ellipsoid (middle); whitening
$\eta = S^{1/2}\alpha$ gives a ball whose maximizer aligns with the
leading eigenvector of $\Gamma$ (right).}
\label{fig:of_ellipsoid}
\end{figure}

\subsection{Saturation Bound}\label{sec:of_sat}

By Remark~\ref{rem:substitution}, the saturation-bound construction
transports verbatim from Section~\ref{sec:rho_sat}.

\begin{proposition}[OF saturation bound]\label{prop:of_rho_sat}
Define
\begin{equation}
\rho_{\mathrm{sat}} = \frac{b}{\bar U_{\max}} - 1.
\label{eq:of_rho_sat}
\end{equation}
If $\rho_{\mathrm{sat}} > 0$ (i.e., $\bar U_{\max} < b$), then for any
$\rho \in [0, \rho_{\mathrm{sat}}]$,
\begin{equation}
\|(1 + \rho)\bar K\,\xi\|_2 \leq b,
\quad \forall\,\xi \in \bar{\mathcal L}_V(\bar c).
\end{equation}
\end{proposition}

\begin{proof}
For any $\xi \in \bar{\mathcal L}_V(\bar c)$,
Lemma~\ref{lemma:Ubar_max} gives $\|\bar K\,\xi\|_2 \leq \bar U_{\max}$.
Then
\begin{align}
\|(1 + \rho)\bar K\,\xi\|_2
&= (1 + \rho)\,\|\bar K\,\xi\|_2 \notag\\
&\leq (1 + \rho)\,\bar U_{\max} \notag\\
&\leq \frac{b}{\bar U_{\max}}\cdot \bar U_{\max} \;=\; b,
\end{align}
where the last inequality uses $\rho \leq \rho_{\mathrm{sat}}$.  At
$\rho = \rho_{\mathrm{sat}}$, equality holds.
\end{proof}

\subsection{Lyapunov Decrease Bound}\label{sec:of_lyap}

The saturation bound alone does not ensure stability.  We bound $\rho$
so that $\bar V(\xi) = \xi^{\!\top}\!\bar P\,\xi$ decreases under
$u_k = (1+\rho)\bar K\,\xi_k$.  Since $\bar H_{uu} = I + B^{\!\top}\!PB$
has the SF form (observation~(ii)) and $\bar H$ enters the Bellman
recursion~\eqref{eq:of_bellman} exactly as $H$ does, the six-step
calculation of Section~\ref{sec:lyapunov} transports verbatim under
$x \mapsto \xi$, $H \mapsto \bar H$, $Kx \mapsto \bar K\xi$; only the
running cost changes, to $\gamma\,y_k^{\!\top}\!y_k$ (still nonnegative).
The result is the OF Lyapunov-decrease identity
\begin{equation}
\bar V(\xi_k) - \bar V(\xi_{k+1})
  = \gamma\,y_k^{\!\top}\!y_k
  + (\bar K\,\xi_k)^{\!\top}\,\bar\Lambda(\rho)\,(\bar K\,\xi_k),
\label{eq:of_lyap_decrease}
\end{equation}
with the OF \emph{Lyapunov-decrease indicator}
\begin{equation}
\boxed{\bar\Lambda(\rho) \;\triangleq\;
       (1+\rho)^2 I_m - \rho^2\,\bar H_{uu},}
\label{eq:Lambda_bar}
\end{equation}
identical in form to the SF indicator~\eqref{eq:Lambda} with
$\bar\lambda := \lambda_{\max}(\bar H_{uu}) > 1$.  The two-quantities
analysis of Section~\ref{sec:lyapunov} therefore applies with $H_{uu}$
replaced by $\bar H_{uu}$---no pseudoinverse enters, since
$\bar H_{uu}$ is a full-rank $m \times m$ block:
\begin{align}
\bar\Lambda(\rho) \succ 0
   \;\iff\; \rho < \rho_{\max}
   &\triangleq \frac{1}{\sqrt{\bar\lambda} - 1},
   \notag\\
\rho_{\mathrm{opt}} \triangleq \frac{1}{\bar\lambda - 1}
   \;&<\; \rho_{\max},
\label{eq:of_rho_opt_derivation}
\end{align}
with $\rho_{\mathrm{opt}}$ the maximizer of the guaranteed decrease
margin, exactly as
in~\eqref{eq:rho_max}--\eqref{eq:rho_order}; the design again selects
$\rho_{\mathrm{opt}}$.

\begin{remark}[More gain-boost room]\label{rem:more_room}
The rank-deficient OF cost (its penalty vanishes on $\ker C$) yields
a ``smaller'' Riccati solution and a $\bar H_{uu}$ closer to the
identity, so both the
maximizer and the stability limit are larger than in SF (Case~1:
$\rho_{\mathrm{opt}} = 4.22$ vs.\ $2.31$ and $\rho_{\max} = 8.91$ vs.\
$5.07$), widening the OF parabola in Fig.~\ref{fig:of_lambda}.  The extra headroom
yields only a modest speedup, however, since the rank-deficient cost
relocates only a subset of the poles (Section~\ref{sec:pole_movement}).
\end{remark}

\begin{figure}[!htb]
\centering
\includegraphics[width=\linewidth]{fig/others (1).pdf}
\caption{State-feedback (dashed) vs.\ output-feedback (solid)
Lyapunov-decrease indicators, Case~1.  Both have the shape
$(1+\rho)^2 - \rho^2\lambda_{\max}(\cdot)$, but the rank-deficient OF
cost penalizes only the output direction, making
$\bar H_{uu} = I + B^{\!\top}\!PB$ closer to the identity; the OF
parabola is therefore wider, enlarging both $\rho_{\mathrm{opt}}$ and
the stability limit $\rho_{\max}$.}
\label{fig:of_lambda}
\end{figure}

\subsection{Stability Theorem}\label{sec:of_stability}

\begin{theorem}[Model-Free OF Stabilization]
\label{thm:of_main}
Under Assumptions~\ref{assum:ANCBC}, \ref{assum:obs},
and~\ref{assum:W_rank}, let $\bar H(\gamma)$ be the converged extended
Q-kernel from the OF policy iteration
of~{\citep[Algorithm~2.1]{RizviLin2023book}}, and let $\bar K$,
$\bar P$, $\bar c$, $\bar U_{\max}$, $\rho_{\mathrm{sat}}$,
$\rho_{\mathrm{opt}}$, and $\rho^{*}$ be computed from $\bar H$ and
$\xi_{\mathrm{init}}$ by the state-feedback formulas of
Theorem~\ref{thm:main} under the single substitution
$P^{-1}\to\bar P^{\dagger}$ (equivalently, Algorithm~\ref{alg:of}).  If
$\rho^{*} > 0$, the boosted controller
$u_k = (1+\rho^{*})\bar K\,\xi_k$ guarantees:
\begin{enumerate}
\item[(a)] \textbf{Saturation avoidance:} $\|u_k\|_\infty \leq b$ for
   all $\xi_k \in \bar{\mathcal L}_V(\bar c)$, so $\sigma(u_k)=u_k$
   along the trajectory;
\item[(b)] \textbf{Asymptotic stability:} the state $x_k = W\,\xi_k$
   converges to zero;
\item[(c)] \textbf{Positive invariance:} $\bar{\mathcal L}_V(\bar c)$
   is positively invariant;
\item[(d)] \textbf{Model-free:} all quantities are computed from the
   learned kernel $\bar H$ and $\xi_{\mathrm{init}}$, with no use of
   $(A, B, C)$ or the parameterization matrix $W$.
\end{enumerate}
Moreover, for any bounded set of initial extended states there exists a
sufficiently small $\gamma > 0$ with $\rho^{*} > 0$ (semi-global
stabilization).
\end{theorem}

\begin{proof}
\emph{(a) Saturation avoidance.}  By Lemma~\ref{lemma:Ubar_max} and
$\rho^{*} \leq \rho_{\mathrm{sat}} = b/\bar U_{\max} - 1$, for any
$\xi_k \in \bar{\mathcal L}_V(\bar c)$,
\begin{equation}
\|u_k\|_\infty \leq \|u_k\|_2
   = (1+\rho^{*})\,\|\bar K\,\xi_k\|_2
   \leq (1+\rho_{\mathrm{sat}})\,\bar U_{\max} = b,
\label{eq:of_sat_chain}
\end{equation}
so $\sigma(u_k) = u_k$ throughout $\bar{\mathcal L}_V(\bar c)$.

\emph{(b) Asymptotic stability.}  With saturation inactive, the OF
decrease identity~\eqref{eq:of_lyap_decrease} holds with
$\bar\Lambda(\rho^{*}) \succ 0$, since
$\rho^{*} \leq \rho_{\mathrm{opt}} < \rho_{\max}$
by~\eqref{eq:of_rho_opt_derivation}; thus $\bar V$ is
non-increasing---all that part~(c) requires.  For convergence, pass to
$x_k = W\xi_k$ (exact for $k \geq N$): by~\eqref{eq:Kbar_collapse},
$\bar K\xi_k = Kx_k$ and $y_k = Cx_k$, while
$\bar V(\xi_k) = x_k^{\!\top}P\,x_k =: V(x_k)$ with $P \succ 0$.  The
closed loop is $x_{k+1} = A_{\mathrm{cl}}\,x_k$ with
$A_{\mathrm{cl}} := A + B(1+\rho^{*})K$,
and~\eqref{eq:of_lyap_decrease} becomes the Lyapunov equation
\begin{equation}
A_{\mathrm{cl}}^{\!\top} P A_{\mathrm{cl}} - P = -S,
\quad
S := \gamma\,C^{\!\top}\!C
     + K^{\!\top}\bar\Lambda(\rho^{*})\,K \;\succeq\; 0.
\label{eq:thm_proof_S}
\end{equation}
Since $\bar\Lambda(\rho^{*}) \succ 0$,
$\ker S = \ker C \cap \ker K$, and $(A_{\mathrm{cl}}, S^{1/2})$ is
detectable: if $A_{\mathrm{cl}}v = \lambda v$ with $S v = 0$, then
$Kv = 0$, so
\begin{equation}
A_{\mathrm{cl}}\,v = \bigl(A + B(1+\rho^{*})K\bigr)\,v = A\,v,
\qquad v \in \ker C,
\label{eq:of_pbh}
\end{equation}
contradicting observability of $(A, C)$
(Assumption~\ref{assum:obs}) by the Popov--Belevitch--Hautus test.
With $P \succ 0$, $S \succeq 0$, and detectability,
\eqref{eq:thm_proof_S} forces $A_{\mathrm{cl}}$ Schur, so
$x_k = W\xi_k \to 0$ within $\bar{\mathcal L}_V(\bar c)$.

\emph{(c) Positive invariance.}  By part~(b),
$\bar V(\xi_{k+1}) \leq \bar V(\xi_k) \leq \bar c$, so
$\xi_{k+1} \in \bar{\mathcal L}_V(\bar c)$.

\emph{(d) Model-free.}  Every quantity in Lines~4--11 of
Algorithm~\ref{alg:of} is a closed-form algebraic function of $\bar H$
and $\xi_{\mathrm{init}}$ alone.

\emph{Semi-global property.}  Exactly as for Theorem~\ref{thm:main},
the closed-loop form of the OF DARE~\eqref{eq:of_dare} is
\begin{equation}
A_{\mathrm{cl}}^{\!\top} P A_{\mathrm{cl}} - P
   = -\bigl(\gamma\,C^{\!\top}\!C + K^{\!\top}\!K\bigr),
\label{eq:of_cl_lyap}
\end{equation}
so $K^{\!\top}\!K \preceq P$ and, with $x = W\xi$,
\begin{equation}
\|\bar K\,\xi\|_2^2 = \|Kx\|_2^2 \leq x^{\!\top}\!P x
   = \xi^{\!\top}\bar P\,\xi \leq \bar c
\quad \text{on } \bar{\mathcal L}_V(\bar c),
\label{eq:of_Umax_bound}
\end{equation}
whence $\bar U_{\max} \leq \sqrt{\bar c}$ and
$\lambda_{\max}(\bar K\,\bar P^{\dagger}\bar K^{\!\top}) \leq 1$
by~\eqref{eq:Ubar_match_sf}.  As $\gamma \to 0$, $P \to 0$, so
$\bar c \to 0$, $\bar U_{\max} \to 0$,
$\rho_{\mathrm{sat}} \to \infty$, and eventually
$\rho^{*} = \rho_{\mathrm{opt}} > 0$, which the $\gamma$-scheduling
loop of Algorithm~\ref{alg:of} locates automatically.
\end{proof}

\subsection{Algorithm}\label{sec:of_algorithm}

Algorithm~\ref{alg:of} summarizes the complete OF procedure.  The
high-gain enhancement (Lines~4--11) consists of closed-form algebraic
operations on the converged extended kernel $\bar H$ and adds
negligible overhead to the underlying output-feedback Q-learning loop.

\begin{algorithm}[!htb]
\caption{Model-Free Low-and-High Gain Output-Feedback Q-Learning}
\label{alg:of}
{\footnotesize
\begin{algorithmic}[1]
\Require $\xi_{\mathrm{init}}$, limit $b$, initial $\gamma$, tolerance $\epsilon$
\Ensure  high-gain OF controller $\bar K_{\mathrm{high}}$
\State Initialize stabilizing $\bar K^{0}$;\ $\gamma_c \leftarrow \gamma$
\Repeat
    \State Run OF PI \citep[Alg.~2.1]{RizviLin2023book} at $\gamma_c$
    \Statex \quad\, until $\|\bar H^{j} - \bar H^{j-1}\| < \epsilon$
    \State $\bar K \leftarrow -\bar H_{uu}^{-1}\,\bar H_{u\xi}$
    \State $\bar P \leftarrow \bar H_{\xi\xi} - \bar H_{\xi u}\,\bar H_{uu}^{-1}\,\bar H_{u\xi}$
    \State $\bar c \leftarrow \xi_{\mathrm{init}}^{\!\top}\,\bar P\,\xi_{\mathrm{init}}$
    \State $\bar U_{\max} \leftarrow \sqrt{\bar c\,\lambda_{\max}(\bar K\,\bar P^{\dagger}\,\bar K^{\!\top})}$
    \State $\rho_{\mathrm{sat}} \leftarrow b/\bar U_{\max} - 1$
    \State $\rho_{\mathrm{opt}} \leftarrow 1/(\lambda_{\max}(\bar H_{uu}) - 1)$
    \State $\rho^{*} \leftarrow \max(0,\min(\rho_{\mathrm{sat}},\rho_{\mathrm{opt}}))$
    \State $\bar K_{\mathrm{high}} \leftarrow (1 + \rho^{*})\,\bar K$
    \If{$\rho_{\mathrm{sat}} \leq 0$}\Comment{$\bar U_{\max}\!\geq\! b$: invariant-set test}
        \State $\gamma_c \leftarrow \gamma_c / 2$
    \EndIf
\Until{$\rho_{\mathrm{sat}} > 0$}
\State \Return $\bar K_{\mathrm{high}}$, $\rho^{*}$, $\gamma_c$
\end{algorithmic}}
\end{algorithm}

The differences from Algorithm~\ref{alg:sf} are mechanical:
input/output data in the policy-iteration step,
$\bar P^{\dagger}$ in place of $P^{-1}$ in the level-set check, and the
same invariant-set $\gamma$-scheduling test.


%% =====================================================================
\section{Numerical Results}\label{sec:results}
%% =====================================================================

The benchmark, taken from~\citep{RizviLin2018CDC,RizviLin2019IJRNC},
is of fourth order with its entire spectrum on the unit circle:
\begin{equation}
A = \begin{bmatrix}
       0 & 1 & 0 & 0 \\
       0 & 0 & 1 & 0 \\
       0 & 0 & 0 & 1 \\
      -1 & 2\sqrt{2} & -4 & 2\sqrt{2}
    \end{bmatrix},
\quad
B = \begin{bmatrix} 0 \\ 0 \\ 0 \\ 1 \end{bmatrix},
\label{eq:system_matrices}
\end{equation}
and $b = 1$.  The eigenvalues $e^{\pm j\pi/4}$, each of multiplicity
two, place the plant at the edge of the ANCBC class.  A trajectory
counts as settled once $\|x_k\| < 0.01\|x_0\|$; two initial conditions
probe the two operating regimes.

\subsection{Case~1: \texorpdfstring{$x_0 = [1,1,1,1]^\top$}{x0 = [1,1,1,1]} (Lyapunov-limited)}

Scheduling terminates at $\gamma = 2.44 \times 10^{-4}$; the low-gain
loop needs 64 steps and never exceeds $\|u\|_{\max} = 0.129$, a mere
12.9\% of the actuator range.  From that very kernel, at that very
$\gamma$, the boost computation returns
\begin{align}
\rho_{\mathrm{sat}} &= 2.88, \quad
\rho_{\mathrm{opt}} = \frac{1}{\lambda_{\max}(H_{uu})-1} = 2.31,
\notag \\
\rho^* &= \min(2.88,\; 2.31) = 2.31, \notag \\
\Lambda(\rho^*) &= (1+2.31)^2 - (2.31)^2 \times 1.433 = 3.31.
\end{align}
Here $\rho_{\mathrm{opt}} < \rho_{\mathrm{sat}}$: \emph{the decrease
certificate, not the actuator, caps the boost}.  Under
$K_{\mathrm{high}} = (1+\rho^{*})K = 3.31\,K$ the loop settles in 37
steps---a $1.73\times$ acceleration---peaking at
$\|u\|_{\max} = 0.318$ (31.8\%) with no violation
(Figs.~\ref{fig:baseline_c1}--\ref{fig:comparison_c1}).

%% ---- CASE 1 FIGURES (grouped: baseline / proposed / overlay) ----
\begin{figure*}[!t]
\centering
\subfloat[Baseline ($\gamma = 2.44\times10^{-4}$): states settle in
64 steps; control uses 12.9\% of capacity.]{%
  \includegraphics[width=0.32\textwidth]{fig/Low-Gain (Original) sf.pdf}%
  \label{fig:baseline_c1}}\hfil
\subfloat[Proposed ($\rho^* = 2.31$, $\Lambda(\rho^*) = 3.31$): 37
steps, 31.8\% utilization; binding $\rho_{\mathrm{opt}}$.]{%
  \includegraphics[width=0.32\textwidth]{fig/High-Gain (Proposed) sf.pdf}%
  \label{fig:proposed_c1}}\hfil
\subfloat[Overlay: baseline (solid) vs.\ proposed (dashed);
$1.73\times$ speedup ($64 \to 37$).]{%
  \includegraphics[width=0.32\textwidth]{fig/Comparison sf.pdf}%
  \label{fig:comparison_c1}}
\caption{State-feedback Case~1 ($x_0 = [1,1,1,1]^\top$,
Lyapunov-limited).}
\label{fig:sf_case1}
\end{figure*}

\subsection{Case~2: \texorpdfstring{$x_0 = [5,5,5,5]^\top$}{x0 = [5,5,5,5]} (Saturation-limited)}

The larger initial condition drives the schedule further, to
$\gamma = 6.10 \times 10^{-5}$; the baseline now takes 92 steps at
$\|u\|_{\max} = 0.477$ (47.7\%).  The boost computation gives
\begin{align}
\rho_{\mathrm{sat}} &= 0.84, \quad
\rho_{\mathrm{opt}} = \frac{1}{\lambda_{\max}(H_{uu})-1} = 3.44,
\notag \\
\rho^* &= \min(0.84,\; 3.44) = 0.84, \notag \\
\Lambda(\rho^*) &= (1+0.84)^2 - (0.84)^2 \times 1.291 = 2.47.
\end{align}
This time $\rho_{\mathrm{sat}} < \rho_{\mathrm{opt}}$: \emph{the
actuator geometry, not the certificate, is limiting}---the mirror
image of Case~1.  Even so, settling drops to 64 steps ($1.44\times$)
at $\|u\|_{\max} = 0.820$ (82.0\%), again without violation
(Figs.~\ref{fig:baseline_c2}--\ref{fig:comparison_c2}).

%% ---- CASE 2 FIGURES (grouped: baseline / proposed / overlay) ----
\begin{figure*}[!t]
\centering
\subfloat[Baseline ($\gamma = 6.10\times10^{-5}$): states settle in
92 steps; control uses 47.7\% of capacity.]{%
  \includegraphics[width=0.32\textwidth]{fig/lg.pdf}%
  \label{fig:baseline_c2}}\hfil
\subfloat[Proposed ($\rho^* = 0.84$, $\Lambda(\rho^*) = 2.47$): 64
steps, 82.0\% utilization; binding $\rho_{\mathrm{sat}}$.]{%
  \includegraphics[width=0.32\textwidth]{fig/hg.pdf}%
  \label{fig:proposed_c2}}\hfil
\subfloat[Overlay: baseline (solid) vs.\ proposed (dashed);
$1.44\times$ speedup ($92 \to 64$).]{%
  \includegraphics[width=0.32\textwidth]{fig/comp.pdf}%
  \label{fig:comparison_c2}}
\caption{State-feedback Case~2 ($x_0 = [5,5,5,5]^\top$,
saturation-limited).}
\label{fig:sf_case2}
\end{figure*}

\subsection{Comparative Discussion}

The two cases visit both branches of $\rho^{*}$
(Table~\ref{tab:sf_vs_of}): headroom in Case~1 lets the certificate
bind ($2.31 < 2.88$), pressure in Case~2 lets the geometry bind
($0.84 < 3.44$), with no intervention either way.  Idle authority is
put to work---$1.73\times$ and $1.44\times$ faster settling,
$\Lambda(\rho^*)\succ 0$, no violation---at the cost of one Schur
complement and one eigenvalue.

%% =====================================================================
%% Output-feedback numerical results
%% =====================================================================

\subsection{Output-Feedback Case~1: Lyapunov-Limited}
\label{sec:of_results_c1}

We now apply the OF framework of Section~\ref{sec:proposed_of} to the
same fourth-order test system~\eqref{eq:system_matrices} with output
matrix $C = [1, 0, 0, 0]$, so that only the first state component is
measured.  The OF cost weighting is
$Q = \gamma\,C^{\!\top}\!C = \gamma\,\mathrm{diag}(1, 0, 0, 0)$, which
has rank one.  The dead-beat parameterization horizon is $N = 4$
(the system order), giving an extended state
$\xi_k \in \mathbb{R}^{(m + p)N} = \mathbb{R}^{8}$ for the SISO case ($m = p = 1$).
The extended Q-kernel
$\bar H \in \mathbb{R}^{9 \times 9}$ has $\bar s = 9 \cdot 10 / 2 = 45$ unique
entries to learn---three times the SF count.

Output-feedback Q-learning {\citep[Algorithm~2.1]{RizviLin2023book}}
converges at $\gamma = 1.22 \times 10^{-4}$, about half the SF value
on the same problem and consistent with the additional headroom
afforded by the rank-1 cost.  The baseline OF controller $\bar K$
settles in 105 steps with peak control $\|u\|_{\max} = 0.096$, i.e.,
$9.6\%$ actuator utilization.  The proposed OF high-gain extension
computes:
\begin{align}
&\bar H_{uu} = 1.237, \quad \bar c = 0.066, \quad \bar U_{\max} = 0.112, \notag\\
&\rho_{\mathrm{sat}} = b/\bar U_{\max} - 1 = 7.93, \notag\\
&\rho_{\mathrm{opt}} = 1/(\bar H_{uu} - 1) = 4.22, \quad \rho^{*} = 4.22, \notag\\
&\bar\Lambda(\rho^{*}) = (1 + 4.22)^2 - (4.22)^2 \times 1.237 = 5.22.
\end{align}
As in SF Case~1, the Lyapunov-decrease bound is binding
($\rho_{\mathrm{opt}} < \rho_{\mathrm{sat}}$).  The boosted controller
$\bar K_{\mathrm{high}} = (1+\rho^{*})\bar K = 5.22\,\bar K$ settles in
85 steps
($1.24\times$ speedup) with $\|u\|_{\max} = 0.370$ ($37.0\%$
utilization) and zero violations
(Figs.~\ref{fig:of_baseline_c1}--\ref{fig:of_comparison_c1}).

Against SF Case~1, the smaller $\bar H_{uu}$ ($1.237$ vs.\ $1.433$)
inflates both bounds ($\rho_{\mathrm{opt}}$ $4.22$ vs.\ $2.31$;
$\rho_{\mathrm{sat}}$ $7.93$ vs.\ $2.88$), yet the realized speedup
shrinks ($1.24\times$ vs.\ $1.73\times$)---the pole-movement effect of
Remark~\ref{rem:more_room}, analyzed in
Section~\ref{sec:pole_movement}.

\begin{figure*}[!t]
\centering
\subfloat[Baseline ($\gamma = 1.22\times10^{-4}$): states settle in
105 steps; control uses 9.6\% of capacity.]{%
  \includegraphics[width=0.32\textwidth]{figof/c1low.pdf}%
  \label{fig:of_baseline_c1}}\hfil
\subfloat[Proposed ($\rho^{*} = 4.22$, $\bar\Lambda(\rho^{*}) = 5.22$):
85 steps, 37.0\% utilization; binding $\rho_{\mathrm{opt}}$.]{%
  \includegraphics[width=0.32\textwidth]{figof/c1high.pdf}%
  \label{fig:of_proposed_c1}}\hfil
\subfloat[Overlay: baseline (solid) vs.\ proposed (dashed);
$1.24\times$ speedup ($105 \to 85$).]{%
  \includegraphics[width=0.32\textwidth]{figof/c1lowhigh.pdf}%
  \label{fig:of_comparison_c1}}
\caption{Output-feedback Case~1 ($x_0 = [1,1,1,1]^\top$,
Lyapunov-limited).}
\label{fig:of_case1}
\end{figure*}

\begin{figure*}[!t]
\centering
\subfloat[Baseline ($\gamma = 1.53\times10^{-5}$): states settle in
184 steps; control uses 27.5\% of capacity.]{%
  \includegraphics[width=0.32\textwidth]{figof/c2low.pdf}%
  \label{fig:of_baseline_c2}}\hfil
\subfloat[Proposed ($\rho^{*} = 2.24$, $\bar\Lambda(\rho^{*}) = 4.81$):
147 steps, 81.8\% utilization; binding $\rho_{\mathrm{sat}}$.]{%
  \includegraphics[width=0.32\textwidth]{figof/c2high.pdf}%
  \label{fig:of_proposed_c2}}\hfil
\subfloat[Overlay: baseline (solid) vs.\ proposed (dashed);
$1.25\times$ speedup ($184 \to 147$).]{%
  \includegraphics[width=0.32\textwidth]{figof/c2lowhigh.pdf}%
  \label{fig:of_comparison_c2}}
\caption{Output-feedback Case~2 ($x_0 = [5,5,5,5]^\top$,
saturation-limited).}
\label{fig:of_case2}
\end{figure*}

\subsection{Output-Feedback Case~2: Saturation-Limited}
\label{sec:of_results_c2}

For the larger initial condition $x_0 = [5, 5, 5, 5]^\top$, the OF
$\gamma$-scheduling converges at $\gamma = 1.53 \times 10^{-5}$, again
about an order of magnitude smaller than the Case~1 value.  The OF
baseline settles in 184 steps with $\|u\|_{\max} = 0.275$ ($27.5\%$
utilization).  The proposed OF high-gain method yields:
\begin{align}
&\bar H_{uu} = 1.133, \quad \bar c = 0.811, \quad \bar U_{\max} = 0.309, \notag\\
&\rho_{\mathrm{sat}} = 2.24, \quad \rho_{\mathrm{opt}} = 7.50, \quad \rho^{*} = 2.24, \notag\\
&\bar\Lambda(\rho^{*}) = (1 + 2.24)^2 - (2.24)^2 \times 1.133 = 4.81.
\end{align}
The saturation geometry is now binding
($\rho_{\mathrm{sat}} < \rho_{\mathrm{opt}}$), exactly matching SF
Case~2's regime.  The boosted controller settles in 147 steps
($1.25\times$ speedup) with $\|u\|_{\max} = 0.818$ ($81.8\%$
utilization) and zero violations
(Figs.~\ref{fig:of_baseline_c2}--\ref{fig:of_comparison_c2}).



\subsection{Output-Feedback Comparative Discussion}
\label{sec:of_comparative}
The OF runs repeat the automatic selection
(Table~\ref{tab:sf_vs_of}): $\rho_{\mathrm{opt}}$ binds in Case~1
($4.22 < 7.93$), $\rho_{\mathrm{sat}}$ in Case~2 ($2.24 < 7.50$),
utilization climbs to $37.0\%$ and $81.8\%$, and the speedups sit
uniformly near $1.24$--$1.25\times$ whichever constraint binds---a
uniformity explained in Section~\ref{sec:pole_movement}.


\subsection{State-Feedback vs.\ Output-Feedback: The Pole-Movement Trade-off}
\label{sec:pole_movement}

\begin{table}[!htb]
\centering
\caption{State-feedback vs.\ output-feedback comparison across both
initial conditions ($b = 1$).  All bold values are model-free
quantities computed from the learned kernel.}
\label{tab:sf_vs_of}
\renewcommand{\arraystretch}{1.25}
\small
\resizebox{\columnwidth}{!}{%
\begin{tabular}{lcccc}
\toprule
& \multicolumn{2}{c}{\textbf{Case~1:} $x_0 = \mathbf{1}_4$}
& \multicolumn{2}{c}{\textbf{Case~2:} $x_0 = 5 \cdot \mathbf{1}_4$} \\
\cmidrule(lr){2-3}\cmidrule(lr){4-5}
\textbf{Metric} & SF & OF & SF & OF \\
\midrule
$\gamma$ converged
   & $2.44\!\times\!10^{-4}$ & $1.22\!\times\!10^{-4}$
   & $6.10\!\times\!10^{-5}$ & $1.53\!\times\!10^{-5}$ \\
$\bar H_{uu}$ (binding eigenvalue)
   & 1.433 & 1.237 & 1.291 & 1.133 \\
\midrule
$\rho_{\mathrm{sat}}$       & 2.88              & 7.93              & 0.84              & 2.24 \\
$\rho_{\mathrm{opt}}$       & 2.31              & 4.22              & 3.44              & 7.50 \\
$\rho^{*}$                  & \textbf{2.31}     & \textbf{4.22}     & \textbf{0.84}     & \textbf{2.24} \\
Binding constraint          & $\rho_{\mathrm{opt}}$ & $\rho_{\mathrm{opt}}$
                            & $\rho_{\mathrm{sat}}$ & $\rho_{\mathrm{sat}}$ \\
\midrule
$\Lambda(\rho^{*})$         & 3.31 & 5.22 & 2.47 & 4.81 \\
\midrule
Settling baseline (steps)   & 64   & 105  & 92   & 184 \\
Settling proposed (steps)   & 37   & 85   & 64   & 147 \\
Speedup                     & $1.73\times$ & $1.24\times$ & $1.44\times$ & $1.25\times$ \\
\midrule
$\|u\|_{\max}$ baseline     & 0.129 & 0.096 & 0.477 & 0.275 \\
$\|u\|_{\max}$ proposed     & 0.318 & 0.370 & 0.820 & 0.818 \\
Utilization baseline        & 12.9\% & 9.6\%  & 47.7\% & 27.5\% \\
Utilization proposed        & 31.8\% & 37.0\% & 82.0\% & 81.8\% \\
Violations                  & 0    & 0    & 0    & 0 \\
\bottomrule
\end{tabular}}
\end{table}

Table~\ref{tab:sf_vs_of} shows a consistent asymmetry: OF takes a
larger $\rho^{*}$ ($4.22$ vs.\ $2.31$ in Case~1) yet delivers a
smaller speedup ($1.24\times$ vs.\ $1.73\times$)---a root-locus effect
of the cost rank.  Scaling the low-gain gain by $(1+\rho)$ drives the poles
of $A + (1+\rho)BK(\gamma)$ toward the zeros of the low-gain loop
$K(\gamma)(zI-A)^{-1}B$, and by Riccati monotonicity the
rank-deficient $Q=\gamma\,C^{\!\top}\!C$ both weakens $P$ (the larger
margin) and fixes the direction of $K(\gamma)$ so that only some loop
zeros sit near the origin.  (Not for lack of penalization---$(A,C)$ observable means the optimal
loop damps every mode; it is the \emph{scaled}, non-optimal gain that
moves only a subset.)  For OF Case~1, the poles of $A + BK$ and
$A + BK_{\mathrm{high}}$ ($K = \bar K W$) are
\begin{align}
\text{Baseline:}\ &\{0.951,\,0.951,\,0.946,\,0.946\},\notag\\
\text{Proposed:}\ &\{0.947,\,0.947,\,\mathbf{0.706},\,\mathbf{0.009}\}:
\end{align}
two poles move sharply, two barely shift, and the slow pair sets
settling; Case~2 is analogous.  The full-rank SF cost, by contrast,
moves all four poles ($1.73\times$).  The rank of $Q$ thus trades
margin against pole-movement effectiveness: full rank, small margin,
uniform damping; rank-deficient, large margin, selective damping.


%% Flush every remaining float here so that no figure or table can be
%% typeset after the start of the Conclusion.
\FloatBarrier

%% =====================================================================
\section{Conclusion}\label{sec:conclusion}
%% =====================================================================

We have combined low-and-high gain feedback with Q-learning into a
fully model-free design for discrete-time saturated systems, in state-
and output-feedback form.  On the state-feedback side, one Schur
complement of the learned kernel exposes the Lyapunov matrix, and from
it both the certified boost and its exact stability limit emerge in
closed form---no identification, no further learning, no search.  The
output-feedback extension builds its kernel from input--output
histories via the parameterization of \citet{RizviLin2023book}; its
central result is that the entire pipeline transports under one
information-preserving substitution of the pseudoinverse for the
inverse of the Lyapunov matrix.  The experiments show automatic selection of the binding constraint,
faster settling without violation, and markedly higher utilization.  Output feedback gives a larger decrease bound but a smaller realized
speedup: the rank of the cost weighting sets a structural trade-off
between gain-boost margin and pole-movement effectiveness, one that to
our knowledge has not been reported in the model-free
saturated-control literature.
Future work includes multi-agent input saturation, learning under
disturbances and measurement noise, and the single-substitution
principle for other rank-deficient parameterizations.


%% =====================================================================
\begingroup\small
\begin{thebibliography}{}

\bibitem[Bertsekas(2019)]{Bertsekas2019}
Bertsekas, D.~P., 2019. Reinforcement Learning and Optimal Control.
Athena Scientific, Belmont, MA.

\bibitem[Dony et al.(2026)]{DonyRizviLin2026CDC}
Dony, M.~N.-A.-A., Rizvi, S.~A.~A., Lin, Z., 2026. Model-free
low-and-high gain Q-learning for discrete-time linear systems with
actuator saturation. In: Proceedings of the IEEE Conference on
Decision and Control. To appear.

\bibitem[Hu and Lin(2001)]{HuLin2001}
Hu, T., Lin, Z., 2001. Control Systems with Actuator Saturation:
Analysis and Design. Birkh\"{a}user, Boston.

\bibitem[Jiang and Jiang(2012)]{JiangJiang2012}
Jiang, Y., Jiang, Z.-P., 2012. Computational adaptive optimal control
for continuous-time linear systems with completely unknown dynamics.
Automatica 48~(10), 2699--2704.

\bibitem[Kleinman(1968)]{Kleinman1968}
Kleinman, D.~L., 1968. On an iterative technique for Riccati equation
computations. IEEE Transactions on Automatic Control 13~(1), 114--115.

\bibitem[Lewis et al.(2012)]{LewisVrabieVamvoudakis2012}
Lewis, F.~L., Vrabie, D., Vamvoudakis, K.~G., 2012. Reinforcement
learning and feedback control: using natural decision methods to design
optimal adaptive controllers. IEEE Control Systems Magazine 32~(6),
76--105.

\bibitem[Lin(1998)]{Lin1998}
Lin, Z., 1998. Low Gain Feedback. Vol. 240 of Lecture Notes in Control
and Information Sciences. Springer, London.

\bibitem[Lin and Saberi(1993)]{LinSaberi1993}
Lin, Z., Saberi, A., 1993. Semi-global exponential stabilization of
linear systems subject to input saturation via linear feedbacks.
Systems \& Control Letters 21~(3), 225--239.

\bibitem[Lin and Saberi(1995)]{LinSaberi1995}
Lin, Z., Saberi, A., 1995. A semi-global low-and-high gain design
technique for linear systems with input saturation. International
Journal of Robust and Nonlinear Control 5~(5), 381--398.

\bibitem[Lin et al.(2000)]{LinSaberiStoorvogelMantri2000}
Lin, Z., Saberi, A., Stoorvogel, A.~A., Mantri, R., 2000. An
improvement to the low gain design for discrete-time linear systems
subject to input saturation---solutions of semi-global output
regulation problems. International Journal of Robust and Nonlinear
Control 10~(3), 117--135.

\bibitem[Rizvi and Lin(2018)]{RizviLin2018CDC}
Rizvi, S.~A.~A., Lin, Z., 2018. Model-free global stabilization of
discrete-time linear systems with saturation via data-driven
Q-learning. In: Proceedings of the IEEE Conference on Decision and
Control. pp. 2635--2640.

\bibitem[Rizvi and Lin(2019)]{RizviLin2019IJRNC}
Rizvi, S.~A.~A., Lin, Z., 2019. An iterative Q-learning scheme for the
global stabilization of discrete-time linear systems subject to
actuator saturation. International Journal of Robust and Nonlinear
Control 29~(9), 2660--2672.

\bibitem[Rizvi and Lin(2023)]{RizviLin2023book}
Rizvi, S.~A.~A., Lin, Z., 2023. Reinforcement Learning-Based Feedback
Control of Linear Systems with Actuator Saturation. Springer, Cham.

\bibitem[Sussmann et al.(1994)]{SussmannSontagYang1994}
Sussmann, H.~J., Sontag, E.~D., Yang, Y., 1994. A general result on the
stabilization of linear systems using bounded controls. IEEE
Transactions on Automatic Control 39~(12), 2411--2425.

\bibitem[Sutton and Barto(2018)]{SuttonBarto2018}
Sutton, R.~S., Barto, A.~G., 2018. Reinforcement Learning: An
Introduction, 2nd Edition. MIT Press, Cambridge, MA.

\bibitem[Wang et al.(2012a)]{WangGripSaberiJohansen2012}
Wang, X., Grip, H.~F., Saberi, A., Johansen, T.~A., 2012a. A new
low-and-high gain feedback design using MPC for global stabilization of
linear systems subject to input saturation. In: Proceedings of the
American Control Conference. pp. 2337--2342.

\bibitem[Wang et al.(2012b)]{WangSaberiStoorvogelSannuti2012}
Wang, X., Saberi, A., Stoorvogel, A.~A., Sannuti, P., 2012b.
Simultaneous global external and internal stabilization of linear
time-invariant discrete-time systems subject to actuator saturation.
Automatica 48~(5), 1100--1104.

\bibitem[Watkins and Dayan(1992)]{WatkinsDayan1992}
Watkins, C.~J.~C.~H., Dayan, P., 1992. Q-learning. Machine Learning
8~(3--4), 279--292.

\bibitem[Yang et al.(1997)]{YangSontagSussmann1997}
Yang, Y., Sontag, E.~D., Sussmann, H.~J., 1997. Global stabilization of
linear discrete-time systems with bounded feedback. Systems \& Control
Letters 30~(1), 1--7.

\end{thebibliography}
\endgroup

\end{document}
